\documentclass[11pt,a4paper]{article}

\usepackage[margin=1in]{geometry}
\usepackage[T1]{fontenc}
\usepackage[utf8]{inputenc}
\usepackage{lmodern}
\usepackage{enumitem}
\usepackage{hyperref}
\usepackage{tabularx}
\usepackage{array}
\usepackage{xcolor}
\usepackage[most]{tcolorbox}
\usepackage{graphicx}

\setlength{\parindent}{0pt}
\setlength{\parskip}{0.55em}
\renewcommand{\arraystretch}{1.15}

\newtcolorbox{tripfitbox}{
  colback=gray!5,
  colframe=black,
  boxrule=0.5pt,
  arc=2mm,
  left=2mm,
  right=2mm,
  top=1mm,
  bottom=1mm,
  width=\textwidth
}

\newcommand{\screenshotplaceholder}[1]{
\begin{center}
\fbox{
\begin{minipage}[c][4.2cm][c]{0.82\textwidth}
\centering
\textit{#1}
\end{minipage}
}
\end{center}
}

\title{CS486 -- Design Brief 4\\Digital Mockup Prototyping}
\author{Tuan Dang Nguyen \& Nicolas Berlin}
\date{}

\begin{document}

\maketitle

\section{Prototype Information}

\textbf{Product name:} TripFit

\textbf{Prototype URL:} \textit{[Insert Figma / FluidUI / prototype URL here]}

\textbf{Video URL:} \textit{[Insert Switchtube / YouTube URL here]}

\textbf{Prototype scope:} The prototype contains \textit{[insert number]} screens and supports three key path scenarios: a group decision flow, a quick duo filtering flow, and a plan-repair flow. It does not include sign-up or registration screens, following the Design Brief 4 instructions.

\section{Vision Statement}

\textit{TripFit helps young independent travelers choose what to do next during short city trips by turning scattered saved ideas and nearby options into a small, transparent shortlist that fits their current mood, time, energy, budget, weather, and group situation.}

\section{Task Tree}

The task tree below follows the structure developed in Design Brief 3. It represents the main user goal and the ordered tasks needed to reach that goal.

\begin{tripfitbox}
\textbf{Root goal: Choose a contextually fitting next activity during a short city trip}
\end{tripfitbox}

\begin{enumerate}[leftmargin=1.5em]
    \item \textbf{Recognize the decision moment}
    \begin{itemize}
        \item Need a next activity
        \item Plan changes
        \item Free time appears
    \end{itemize}

    \item \textbf{Express the current context}
    \begin{itemize}
        \item Practical context: time left, current location, budget, weather
        \item Experiential context: energy level, mood, desired atmosphere
        \item Social context: solo / duo / group, group preferences
    \end{itemize}

    \item \textbf{Gather candidate options}
    \begin{itemize}
        \item Reuse known options: saved ideas, existing plan, backup ideas
        \item Find realistic local options: nearby places, currently feasible places
    \end{itemize}

    \item \textbf{Filter and compare}
    \begin{itemize}
        \item Remove unrealistic options: too far, closed, too expensive, poor fit
        \item Create a shortlist: few strong options and one backup option
        \item Explain fit: why each option fits now and what trade-offs it involves
    \end{itemize}

    \item \textbf{Choose and act}
    \begin{itemize}
        \item Coordinate the choice
        \item Select the best option
        \item Go to the chosen place or switch to backup if needed
        \item Finish with confidence
    \end{itemize}
\end{enumerate}

\section{Task Tree New}

The task tree below follows the structure developed in Design Brief 3. It represents the main user goal, the ordered subtasks, and the main alternative paths users may follow to reach that goal.

\begin{tripfitbox}
\textbf{Root goal: Choose a contextually fitting next activity during a short city trip}
\end{tripfitbox}

\begin{enumerate}[leftmargin=1.5em]
    \item \textbf{Start a travel decision}
    \begin{itemize}
        \item Decide what to do now
        \item \textbf{OR} find something nearby quickly
        \item \textbf{OR} fix a changed plan
    \end{itemize}

    \item \textbf{Express the current context}
    \begin{itemize}
        \item Practical context: time left, current location, budget, weather
        \item Experiential context: energy level, mood, desired atmosphere
        \item Social context: solo / duo / group, group preferences
    \end{itemize}

    \item \textbf{Gather candidate options}
    \begin{itemize}
        \item Reuse known options: saved ideas, existing plan, backup ideas
        \item \textbf{OR} find realistic local options: nearby places, currently feasible places
    \end{itemize}

    \item \textbf{Filter and compare options}
    \begin{itemize}
        \item Remove unrealistic options: too far, closed, too expensive, poor fit
        \item Create a shortlist: few strong options and one backup option
        \item Explain fit: why each option fits now and what trade-offs it involves
    \end{itemize}

    \item \textbf{Choose and act}
    \begin{itemize}
        \item Coordinate the choice with the group if needed
        \item Select the best option
        \item Save a backup option
        \item Open directions or update the plan
    \end{itemize}
\end{enumerate}



\section{Key Path Scenarios Used for the Prototype}

The prototype was built around three key path scenarios derived from the context scenarios in Design Brief 3.

\subsection{KPS 1: Group decision flow -- incoherence with the app}

Lina is in Barcelona with three friends. They have three hours before dinner, low-to-medium energy, a moderate budget, and many saved ideas from TikTok, Instagram, screenshots, and Google Maps. She opens TripFit, chooses ``Decide what to do now'', confirms the current location and time left, selects the group context, imports saved ideas, and receives a shortlist of three options. She reviews the explanations, compares trade-offs, uses a quick group alignment step, chooses one option, and saves a backup.

\subsection{KPS 2: Quick duo filtering flow}

Tran is in Lyon with one friend and has two hours left before taking the train. She wants something close, low-effort, and practical. She opens TripFit, chooses ``Quick nearby filter'', enters time, energy, and preference information, receives three nearby options ranked by travel time, effort, opening status, and fit, then chooses one and confirms the next action.

\subsection{KPS 3: Plan repair flow -- incoherence with the app}

Claire is in Copenhagen and her planned outdoor activity becomes unrealistic because of rain. She opens TripFit, chooses ``Fix changed plan'', selects ``weather changed'', sees the current gap in her plan, receives three low-friction backup options, chooses a replacement, and confirms the updated plan.

\section{Application of Interaction Design Principles}

This section explains how the prototype applies eight interaction design principles required in Design Brief 4.

\subsection{Clear visibility of user actions in the UI}

TripFit keeps the user’s selected context visible before showing recommendations. For example, in the group decision flow, the context summary displays the selected time left, energy level, budget, weather, group size, and desired atmosphere. This helps users understand what information the system is using and reduces uncertainty about the current state of the interaction.

\textit{Example screen(s):} \textit{[Insert screen number/name, e.g., Context Summary screen]}

\screenshotplaceholder{Optional screenshot: context summary screen showing selected time, mood, energy, budget, weather, and group size.}

\subsection{Design constraints and visual mapping between tasks and controls}

The prototype constrains user input through simple choices such as context chips, cards, and predefined decision modes. Instead of forcing users to describe their situation in a long text field, TripFit maps common travel decision needs to clear actions such as ``Decide what to do now'', ``Quick nearby filter'', and ``Fix changed plan''. This reduces the gulf of execution because users can immediately identify which path matches their goal.

\textit{Example screen(s):} \textit{[Insert screen number/name, e.g., Start Decision screen]}

\screenshotplaceholder{Optional screenshot: start screen with the three decision modes.}

\subsection{Meaningful feedback to ease evaluation}

TripFit gives feedback through the shortlist explanations. Each recommendation includes reasons such as distance, price level, opening status, atmosphere, energy fit, group fit, and trust signals. This helps users evaluate the system’s output instead of receiving a black-box recommendation. The explanation ``why this fits now'' directly addresses the trust issue identified during the interviews.

\textit{Example screen(s):} \textit{[Insert screen number/name, e.g., Shortlist screen or Option Detail screen]}

\screenshotplaceholder{Optional screenshot: recommendation card with ``Why this fits now'' explanation.}

\subsection{Breaking a big task into smaller steps}

Choosing what to do next during a short trip can be cognitively heavy because users must consider time, distance, budget, mood, energy, weather, and group preferences. TripFit breaks this large decision into smaller steps: choosing a decision mode, expressing current context, gathering options, viewing a shortlist, comparing explanations, and confirming a choice. This supports piece-wise problem solving and avoids overwhelming users with all information at once.

\textit{Example screen(s):} \textit{[Insert screen numbers/names showing the sequence]}

\subsection{Task closure}

After the user selects an option, TripFit provides a clear confirmation screen. The screen confirms the chosen place, reminds the user that a backup option has been saved, and presents the next action, such as opening directions or sharing with the group. This creates task closure: the user knows that the decision is complete and what to do next.

\textit{Example screen(s):} \textit{[Insert screen number/name, e.g., Final Choice Confirmation screen]}

\screenshotplaceholder{Optional screenshot: final confirmation screen with chosen option and backup saved.}

\subsection{Congratulating the user}

The final confirmation screen uses positive feedback such as ``Good choice -- backup saved'' or ``This fits your time, budget, and group mood''. This is not only decorative; it reinforces the user’s confidence and reduces the anxiety that the choice may be wrong. This is especially important because the interviews showed that users often hesitate even after finding many options.

\textit{Example screen(s):} \textit{[Insert screen number/name]}

\subsection{Clear exit marks}

The prototype includes clear ways to leave or redirect the flow. Users can go back, edit context, return to the shortlist, or choose a different decision mode. This matters because travel decisions are unstable: weather, group mood, opening hours, or energy can change quickly. Clear exit marks prevent the user from feeling trapped in one recommendation path.

\textit{Example screen(s):} \textit{[Insert screen number/name, e.g., Shortlist screen with Back/Edit Context]}

\subsection{Allowing errors and emergency exits}

TripFit allows users to revise inputs, change context, reject a recommendation, and switch to a backup option. This is important because the system may misunderstand the user’s mood or because the situation may change after the recommendation is shown. The backup option acts as an emergency exit: if the chosen place is too crowded, closed, or no longer appealing, the user can recover without restarting the whole decision process.

\textit{Example screen(s):} \textit{[Insert screen number/name, e.g., Backup Option screen or Edit Context screen]}

\section{Visual Design Principles}

The prototype uses a consistent mobile-first visual style across all flows. The main design decisions are:

\begin{itemize}
    \item \textbf{Visual hierarchy:} main actions such as choosing a decision mode, generating the shortlist, and confirming the final choice are visually prioritized.
    \item \textbf{Grouping:} related information is grouped into cards, such as context summary, recommendation cards, and explanation blocks.
    \item \textbf{Alignment and spacing:} screens use consistent spacing and alignment to make the interface predictable.
    \item \textbf{Color and emphasis:} color is used to highlight selected context, recommended options, and confirmation states.
    \item \textbf{Screen inertia:} the landing screen and subsequent screens keep a similar look and feel, so users do not feel that they are moving between unrelated systems.
\end{itemize}

\textit{To complete after prototype review: Add one or two sentences about the actual color palette, typography, and card style used in your Figma prototype.}

\section{Peer Evaluation and Iteration}

We prepared the following task scenarios for peer walkthroughs.

\subsection{Walkthrough Task 1: Group decision}

You are in Barcelona with three friends. You have three hours before dinner, everyone is slightly tired, and you have many saved places from TikTok, Instagram, screenshots, and Google Maps. Use TripFit to choose what to do next and save a backup option.

\subsection{Walkthrough Task 2: Quick nearby filter}

You are in Lyon with one friend and have two hours before your train. You want a nearby, low-effort activity that still feels worth the time. Use TripFit to find and choose one option.

\subsection{Walkthrough Task 3: Plan repair}

You are in Copenhagen and your planned outdoor activity is no longer realistic because of rain. Use TripFit to find a replacement activity that fits the remaining time before dinner.

\subsection{Observed problems}

\textit{Fill this section after peer walkthroughs.}

\begin{itemize}
    \item \textbf{Problem 1:} [Describe the first issue observed during walkthrough.]
    \item \textbf{Problem 2:} [Describe the second issue observed during walkthrough.]
    \item \textbf{Problem 3:} [Describe the third issue observed during walkthrough.]
\end{itemize}

\subsection{Changes made after feedback}

\textit{Fill this section after improving the prototype.}

\begin{itemize}
    \item \textbf{Change 1:} [Explain what was changed and which principle it addresses.]
    \item \textbf{Change 2:} [Explain what was changed and which principle it addresses.]
    \item \textbf{Change 3:} [Explain what was changed and which principle it addresses.]
\end{itemize}

\section{Conclusion}

TripFit’s prototype demonstrates the central design direction developed across the previous design briefs: the product should not generate more travel content, but help users make a confident decision from the content and options they already have. The prototype focuses on reducing decision overload, supporting contextual fit, explaining recommendations transparently, and helping users finish the decision with a clear next action.

\end{document}